\documentclass[11pt]{article}

\usepackage[a4paper,margin=1in]{geometry}
\usepackage{amsmath,amssymb}
\usepackage{enumitem}
\usepackage{setspace}
\usepackage{hyperref}

\setstretch{1.1}

\begin{document}

\begin{center}
\textbf{\Large CSE 400: Fundamentals of Probability in Computing} \\[4pt]
\textbf{School of Engineering and Applied Sciences, Ahmedabad University} \\[8pt]
\textbf{\Large CSE 400 – Tutorial 2} \\[4pt]
\textbf{Exam Preparation Lecture Scribe}
\end{center}

\hrule
\vspace{0.5cm}

\section*{Q1. Random Point on a Line Segment}

\subsection*{Problem}

A point is chosen at random on a line segment of length $L$. Interpret this statement, and find the probability that the ratio of the shorter to the longer segment is less than $\frac{1}{4}$.

\subsection*{Concept Used}

\begin{itemize}
\item Uniform distribution on an interval
\end{itemize}

\subsection*{Step-by-Step Reasoning}

Let the random point be at distance $X$ from one endpoint.

Since the point is chosen at random on a segment of length $L$:
\[
X \sim \text{Uniform}(0, L)
\]

The two segments formed are:
\[
X \quad \text{and} \quad L - X
\]

The shorter segment is:
\[
\min(X, L-X)
\]

The longer segment is:
\[
\max(X, L-X)
\]

We require:
\[
\frac{\min(X, L-X)}{\max(X, L-X)} < \frac{1}{4}
\]

\subsubsection*{Case 1: $0 < X \le \frac{L}{2}$}

Then:
\[
\min(X, L-X) = X
\]
\[
\max(X, L-X) = L-X
\]

Condition becomes:
\[
\frac{X}{L-X} < \frac{1}{4}
\]

Multiply both sides by $(L-X)$:
\[
X < \frac{1}{4}(L-X)
\]
\[
X < \frac{L}{4} - \frac{X}{4}
\]
\[
X + \frac{X}{4} < \frac{L}{4}
\]
\[
\frac{5X}{4} < \frac{L}{4}
\]
\[
5X < L
\]
\[
X < \frac{L}{5}
\]

\subsubsection*{Case 2: $\frac{L}{2} < X < L$}

Then:
\[
\min(X, L-X) = L-X
\]
\[
\max(X, L-X) = X
\]

Condition becomes:
\[
\frac{L-X}{X} < \frac{1}{4}
\]
\[
L-X < \frac{X}{4}
\]
\[
L < \frac{5X}{4}
\]
\[
4L < 5X
\]
\[
X > \frac{4L}{5}
\]

\subsection*{Favourable Region}

\[
X < \frac{L}{5}
\]
or
\[
X > \frac{4L}{5}
\]

Total favourable length:
\[
\frac{L}{5} + \frac{L}{5} = \frac{2L}{5}
\]

Total possible length = $L$

\subsection*{Probability}

\[
P = \frac{\frac{2L}{5}}{L}
\]
\[
P = \frac{2}{5}
\]

\newpage

\section*{Q2. Normal Distributions and Error Probability}

\subsection*{Concepts Used}

\begin{itemize}
\item Normal distribution
\item Bayes’ classification principle
\item Conditional probability
\end{itemize}

White region:
\[
X \sim N(4,4)
\]

Black region:
\[
X \sim N(6,9)
\]

Fraction black = $\alpha$

Given reading = 5

Error probability equal means:
\[
P(\text{white}|5) = P(\text{black}|5)
\]

Using Bayes’ rule:
\[
\alpha f_B(5) = (1-\alpha) f_W(5)
\]

Compute densities:
\[
f_W(5) = \frac{1}{\sqrt{2\pi(4)}} e^{-\frac{(5-4)^2}{2(4)}}
\]
\[
= \frac{1}{2\sqrt{2\pi}} e^{-\frac{1}{8}}
\]

\[
f_B(5) = \frac{1}{\sqrt{2\pi(9)}} e^{-\frac{(5-6)^2}{18}}
\]
\[
= \frac{1}{3\sqrt{2\pi}} e^{-\frac{1}{18}}
\]

Thus:
\[
\alpha \frac{1}{3\sqrt{2\pi}} e^{-1/18}
=
(1-\alpha) \frac{1}{2\sqrt{2\pi}} e^{-1/8}
\]

Solve for $\alpha$.

\section*{Q3. Minimizing $E|X-a|$}

\subsection*{(a) Uniform on $(0,A)$}

Concept: Uniform distribution

\[
X \sim U(0,A)
\]

\[
E|X-a| = \int_0^A |x-a|\frac{1}{A} dx
\]

Split at $a$:
\[
= \frac{1}{A}\left[
\int_0^a (a-x)dx + \int_a^A (x-a)dx
\right]
\]

Compute each integral step by step.

Minimization yields:
\[
a = \frac{A}{2}
\]

\subsection*{(b) Exponential with rate $\lambda$}

Concept: Exponential distribution

\[
f(x)=\lambda e^{-\lambda x}
\]

\[
E|X-a|=\int_0^\infty |x-a|\lambda e^{-\lambda x}dx
\]

Split at $a$.

Minimization yields:
\[
a = \frac{\ln 2}{\lambda}
\]

\section*{Q4. Exponential Mixture}

Concepts:
\begin{itemize}
\item Exponential distribution
\item Conditional probability
\end{itemize}

\[
P(T>t+s|T>t)
=
\frac{P(T>t+s)}{P(T>t)}
\]

Mixture survival function:
\[
P(T>t)
=
p_1 e^{-\lambda_1 t}
+
p_2 e^{-\lambda_2 t}
\]

Thus:
\[
\frac{
p_1 e^{-\lambda_1 (t+s)}
+
p_2 e^{-\lambda_2 (t+s)}
}{
p_1 e^{-\lambda_1 t}
+
p_2 e^{-\lambda_2 t}
}
\]

\section*{Q5. Independent Poisson Variables}

Concept:
\begin{itemize}
\item Poisson distribution
\item Independence
\end{itemize}

\[
P(X=k|X+Y=n)
=
\frac{P(X=k, Y=n-k)}{P(X+Y=n)}
\]

Using Poisson formula:
\[
\frac{
e^{-\lambda_1}\frac{\lambda_1^k}{k!}
\cdot
e^{-\lambda_2}\frac{\lambda_2^{n-k}}{(n-k)!}
}{
e^{-(\lambda_1+\lambda_2)}\frac{(\lambda_1+\lambda_2)^n}{n!}
}
\]

Simplify:
\[
\binom{n}{k}
\left(
\frac{\lambda_1}{\lambda_1+\lambda_2}
\right)^k
\left(
\frac{\lambda_2}{\lambda_1+\lambda_2}
\right)^{n-k}
\]

\section*{Q6. Transformation $Y=g(X)$}

Concept:
\begin{itemize}
\item Uniform distribution
\item CDF definition
\item PDF as derivative of CDF
\end{itemize}

\[
F_Y(y)=P(Y\le y)
\]

Compute using graphical mapping of intervals.

Then:
\[
f_Y(y)=\frac{d}{dy}F_Y(y)
\]

\section*{Q7. Stock Movement Model}

Concept:
\begin{itemize}
\item Independent Bernoulli trials
\item Binomial approximation
\end{itemize}

After 1000 periods:
\[
\text{Price}=s u^k d^{1000-k}
\]

Condition:
\[
u^k d^{1000-k} \ge 1.3
\]

Take log:
\[
k\ln u + (1000-k)\ln d \ge \ln 1.3
\]

Solve for $k$.

Probability:
\[
P(K\ge k_0)
=
\sum_{k=k_0}^{1000}
\binom{1000}{k}p^k(1-p)^{1000-k}
\]

\section*{Q8. Exponential Distribution}

\[
\lambda=\frac{1}{2}
\]

\subsection*{(a)}

\[
P(T>2)=e^{-\lambda 2}
=
e^{-1}
\]

\subsection*{(b)}

\[
P(T>10|T>9)
=
\frac{e^{-\lambda 10}}{e^{-\lambda 9}}
=
e^{-\lambda}
=
e^{-1/2}
\]

\section*{Q9. Distribution of $Z=\sqrt{X^2+Y^2}$}

Concept:
\begin{itemize}
\item Joint pdf of independent variables
\item Transformation
\end{itemize}

Joint pdf:
\[
f_{X,Y}(x,y)=f_X(x)f_Y(y)
\]

\[
F_Z(z)=P(X^2+Y^2 \le z^2)
\]

\[
\int\int_{x^2+y^2\le z^2}
f_X(x)f_Y(y)\,dx\,dy
\]

Differentiate:
\[
f_Z(z)
=
\frac{d}{dz}F_Z(z)
\]

\section*{Q10. Joint Distribution of $M_n$ and $M_{n+1}$}

Concept:
\begin{itemize}
\item Maximum of i.i.d. random variables
\end{itemize}

\[
M_n=\max(X_1,\dots,X_n)
\]

\[
P(M_n\le x)
=
F(x)^n
\]

\[
M_{n+1}=\max(M_n, X_{n+1})
\]

Joint distribution:
\[
P(M_n\le x, M_{n+1}\le y)
\]

Case analysis:

If $y<x$: probability 0.

If $y\ge x$:
\[
P(M_n\le x, X_{n+1}\le y)
=
F(x)^n F(y)
\]

\vspace{0.5cm}
\hrule
\vspace{0.3cm}

\begin{center}
\textbf{End of Tutorial-2 Exam Scribe}
\end{center}

\end{document}